\chapter{Total Knee Replacement}
\section*{Overview}
Total knee replacement removes damaged joint surfaces and replaces them with metal and plastic components to relieve pain and restore function.
\section*{Why It Is Done}
It is commonly performed for severe osteoarthritis, inflammatory arthritis, or post-traumatic joint damage causing persistent pain and disability.
\section*{How It Is Performed}
Under regional or general anesthesia, damaged portions of the femur, tibia, and sometimes patella are resurfaced and the components are aligned and secured.
\section*{Risks and Recovery}
Infection, blood clots, stiffness, fracture, nerve injury, instability, loosening, and revision surgery are possible. Early mobilization and structured rehabilitation are essential.
\section*{References}
\begin{enumerate}\item MedlinePlus. Knee Replacement. U.S. National Library of Medicine; 2025.\item Current Medical Diagnosis and Treatment. McGraw-Hill; 2025.\end{enumerate}
\clearpage
